% Fortsetzung von Kapitel 2. Quelle: kumulative Mitschrift, PDF-S. 27--33.
\subsection{Die Russellsche Antinomie}
\notizquelle{27--28}

Mengen können selbst Elemente anderer Mengen sein. Dabei sind
Element- und Teilmengenrelation sorgfältig zu unterscheiden:
$\{0\}\ne\{\{0\}\}$, und
$\{\{0\},0,\{0,1\}\}$ enthält drei verschiedene Elemente.
Für jede Menge $M$ gilt $M\subset M$; daraus folgt keineswegs $M\in M$.

\begin{satz}[Russellsche Antinomie]
Die uneingeschränkte Forderung, zu jeder Eigenschaft eine Menge aller
Objekte mit dieser Eigenschaft bilden zu dürfen, führt zu einem Widerspruch.
\end{satz}
\begin{proof}
Angenommen, die Zusammenfassung
\[
    R:=\{M\mid M\text{ ist eine Menge und }M\notin M\}
\]
wäre eine Menge. Die definierende Eigenschaft müsste dann auch für $R$
selbst gelten, also
\[
    R\in R\quad\Longleftrightarrow\quad R\notin R.
\]
Ist $R\in R$, so folgt $R\notin R$; ist $R\notin R$, so folgt $R\in R$.
Beide Fälle widersprechen der Annahme.
\end{proof}

\begin{beispiel}[Der Barbier]
Ein Barbier soll genau die Personen rasieren, die sich nicht selbst rasieren.
Gehört er selbst zur betrachteten Personengruppe, so müsste er sich genau
dann selbst rasieren, wenn er sich nicht selbst rasiert. Ein Barbier mit
dieser Eigenschaft kann nicht existieren. Dies ist eine Veranschaulichung
der Antinomie, kein Beweis über tatsächlich existierende Personen.
\end{beispiel}

\begin{bemerkung}[Lügner und Epimenides]
Die Mitschrift nennt auch den Satz eines Kreters: \glqq Alle Kreter
lügen\grqq. Dieser ist ohne weitere Annahmen kein logisches Paradoxon:
Ist er falsch, folgt nur, dass nicht alle Kreter immer lügen.
Der eigentliche Lügner-Satz \glqq Dieser Satz ist falsch\grqq\ richtet
sich dagegen unmittelbar auf seinen eigenen Wahrheitswert.
\end{bemerkung}

\subsection{Axiomatik und Klassen}
\notizquelle{28--32}

Die Mitschrift skizziert verschiedene Reaktionen auf solche Antinomien:
Typentheorien trennen verschiedene Stufen von Objekten; eine Beschränkung
auf endliche Mengen wäre für die hier benötigte Analysis zu eng.
Der in der Vorlesung verfolgte Weg ist die \textbf{Axiomatik}.

\begin{definition}[Axiomatischer Ansatz]
Axiome sind festgelegte Grundannahmen einer Theorie. Ein Beweis leitet
weitere Aussagen durch logische Schlussregeln aus diesen Annahmen ab.
Axiome müssen nicht als unmittelbar \glqq offensichtlich wahr\grqq\ gelten;
entscheidend ist, welche Theorie sie beschreiben.
\end{definition}

Für die Mengenlehre verwenden wir den Rahmen von
\textbf{Zermelo--Fraenkel (ZF)}. Mit dem zusätzlichen Auswahlaxiom spricht
man von \textbf{ZFC}. Die Vorlesung gibt hier einen Überblick, keine
vollständige Axiomenliste.

\begin{description}[style=nextline,leftmargin=0pt]
\item[Extensionalitätsaxiom]
Zwei Mengen sind genau dann gleich, wenn sie dieselben Elemente besitzen:
\[
    M=N\quad\Longleftrightarrow\quad
    \forall x:\ (x\in M\Longleftrightarrow x\in N).
\]
Insbesondere gibt es höchstens eine leere Menge. Ihre Existenz folgt aus
den Existenzaxiomen von ZF; Extensionalität allein liefert nur Eindeutigkeit.
\item[Aussonderungsschema]
Ist $M$ bereits eine Menge und $P(x)$ eine zulässige Eigenschaft, dann ist
\[
    \{x\in M\mid P(x)\}
\]
ebenfalls eine Menge. Die Einschränkung auf eine schon vorhandene Menge
ist wesentlich: Eine uneingeschränkte Menge $\{x\mid P(x)\}$ wird nicht
zugelassen.
\item[Weitere Existenzforderungen]
Aus den entsprechenden ZF-Axiomen folgt: Für Mengen $M,N$ sind $M\cup N$
und $\Pot(M)$ Mengen, und es existiert eine unendliche Menge.
Das Auswahlaxiom ist eine zusätzliche Annahme und wird hier nicht
weiter ausgeführt.
\end{description}

\begin{beispiel}[Aussonderung in einer gegebenen Menge]
Der Schnitt ist durch Aussonderung aus $M$ konstruierbar:
\[
    M\cap N=\{x\in M\mid x\in N\}.
\]
Ebenso ist $M\setminus N=\{x\in M\mid x\notin N\}$ eine Menge.
\end{beispiel}

\begin{bemerkung}[Mengen und echte Klassen]
Eine durch eine Eigenschaft beschriebene Gesamtheit von Mengen nennen wir
hier eine \textbf{Klasse}. In ZF ist dies eine abkürzende Sprechweise,
keine zusätzliche Art von Objekten innerhalb der Theorie.
Jede Menge bestimmt eine Klasse, aber nicht jede Klasse ist eine Menge.
Eine Klasse, die keine Menge ist, heißt \textbf{echte Klasse}.
Die Gesamtheit aller Mengen ist eine echte Klasse: Wäre sie eine Menge
$U$, könnte man durch Aussonderung
$\{M\in U\mid M\notin M\}$ bilden und erhielte erneut Russells Widerspruch.
\end{bemerkung}

\subsection{Unvollständigkeit}
\notizquelle{32--33}

Die Mitschrift gibt einen informellen Ausblick auf Gödels
Unvollständigkeitssätze (1931). Die dortige Formulierung
\glqq jede halbwegs interessante Theorie\grqq\ benötigt folgende
Voraussetzungen.

\begin{satz}[Erster Unvollständigkeitssatz, ohne Beweis]
Eine widerspruchsfreie, effektiv axiomatisierbare formale Theorie, die
genügend elementare Arithmetik ausdrücken kann, ist unvollständig:
Es gibt eine Aussage $A$ ihrer Sprache, für die weder $A$ noch $\neg A$
in dieser Theorie beweisbar ist.
\end{satz}

\begin{satz}[Zweiter Unvollständigkeitssatz, ohne Beweis]
Eine widerspruchsfreie, effektiv axiomatisierbare Theorie mit hinreichend
viel Arithmetik kann ihre eigene, in üblicher Weise formalisierte
Widerspruchsfreiheit nicht innerhalb derselben Theorie beweisen.
\end{satz}

\begin{bemerkung}[Einordnung]
\glqq Unvollständig\grqq\ bedeutet nicht, dass eine Aussage weder wahr
noch falsch ist, sondern dass die gewählten Axiome sie nicht entscheiden.
Ein Beweis der Widerspruchsfreiheit in einer stärkeren Theorie ist damit
nicht ausgeschlossen. Die Sätze sind auch kein Nachweis, dass ZF
widersprüchlich ist. Das Buch-Beispiel der Mitschrift dient lediglich
der Anschauung und ersetzt keinen Beweis.
\end{bemerkung}

\noindent{\small Fachliche Präzisierung der Voraussetzungen: P.~Raatikainen,
\emph{Gödel's Incompleteness Theorems}, Stanford Encyclopedia of Philosophy,
\url{https://plato.stanford.edu/entries/goedel-incompleteness/}.}
